\begin{ejercicio}
    Sea $A$ un anillo. Diremos que $A$ es trivial si $A=\{0\}$. Demostrar que $A$ es trivial si, y solo si, $1=0$. \\

\noindent
\textbf{Solución:}
\begin{description}
    \item [$\Longrightarrow )$] Es claro, ya que $A = \{0\}$. 
    \item [$\Longleftarrow )$] Sea $a \in A$, tenemos que $a = a\cdot 1 = a \cdot 0 = 0 \Longrightarrow a = 0$. Por tanto, $A$ es trivial.
\end{description}
\end{ejercicio}

%--------------------------------------------------------------
\begin{ejercicio}
    Sea $K$ un cuerpo y $M_n(K)$ el anillo de matrices cuadradas de orden $n$ con entradas en $K$. Demostrar que $Z(M_n(K)) = \{kI_n | k \in K\}$, donde $I_n$ es la matriz identidad de orden $n$. \\

\noindent
\textbf{Solución:}

Está claro que el conjunto definido está dentro del centro. Sea $A \in M_n(K)$, $k \in K$:
\begin{equation*}
    kI_n A = kA = kAI_n = AkI_n
\end{equation*}
Ahora debemos ver que si $A \in Z(M_n(K)) \Longrightarrow A = kI_n$.
Defino la matriz $E_{ml} = (e_{ij}) \in M_n(K)$ tal que $e_{ij} = 1$ si $i = m$ y $j = l$ y $e_{ij} = 0$ en otro caso.
Sea ahora $A \in M_n(K)$. Vamos a igualar $E_{ml}A = AE_{ml}$ y vamos a sacar condiciones para la $A$. Veremos que si tomamos $m,l \in \{1\ldots n\}$, $A= (a_{ij})$ deberá ser de la forma $kI_n$ con $k \in K$. \\

\noindent
Sea $(c_{ij}) = E_{ml}A$.
\begin{equation*}
    c_{ij} = \sum_{t=1}^n e_{it}a_{tj} =  \left\{\begin{array}{ll}
        0 & \text{si\ } i \neq m \\
        a_{lj} & \text{si\ i = m} 
    \end{array}\right. 
\end{equation*}
Sea ahora $(d_{ij}) = AE_{ml}$.
\begin{equation*}
    d_{ij} = \sum_{t=1}^n a_{it}e_{tj} =  \left\{\begin{array}{ll}
        0 & \text{si\ } j \neq k \\
        a_{im} & \text{si\ j = k} 
    \end{array}\right. 
\end{equation*}
Ahora, si forzamos a que sean iguales tenemos los siguientes casos:
\begin{gather*}
    a_{lj} = 0 \text{\ si\ } i=m \land j\neq l\\
    a_{im} = 0 \text{\ si\ } i\neq m \land j=l \\
    a_{lj} = a_{im} \text{\ si\ } i=m \land j=l \text{\ es\ decir\ } a_{ll} = a_{mm} \\ 
    0 = 0 \text{\ si\ } i\neq m \land j\neq l
\end{gather*}

De esta forma vemos que todos los elementos de la diagonal deben ser iguales y todos los que no son de la diagonal deben anularse, es decir:
\begin{equation*}
    A = kI_n \text{\ con\ } k \in K
\end{equation*}

\end{ejercicio}

%--------------------------------------------------------------
\begin{ejercicio}
    Sea $V$ un espacio vectorial sobre un cuerpo $K$. El conjunto
    \begin{equation*}
        End_k(V) =\{f:V\rightarrow V : f \text{\ es\ } K-lineal\}
    \end{equation*}
    es un subanillo de $End(V)$. Consideramos la aplicación $h:K \rightarrow End_k(V)$ que asigna a cada $k \in K$ la homotecia $h(k):V \rightarrow V$, definido por $h(k)(v) = kv$ $\forall v \in V$. Comprobar que $h$ está bien definida y que es un homomorfismo de anillos. 

    Además, si $T:V \rightarrow V$ $K$-lineal y $k \in K$, tenemos que $T \circ h(k) = h(k) \circ T$, luego $Im(h) \subset Z(End_k(V))$. Y así $End_K(V)$ es una $K$-álgebra.

\noindent
\textbf{Solución:}
\begin{itemize}
    \item Para ver que es un subanillo lo que falta ver es que las operaciones sacan los escalares, sean $f, g \in End_K(V)$
        \begin{gather*}
            (f-g)(kv) = f(kv) - g(kv) = kf(v) - kg(v) = k (f-g)(v). \\
            (f \circ g) (kv) = f(g(kv)) = f(kg(v))= kf(g(v)) = k(f\circ g)(v)
        \end{gather*}

    \item Para ver que $h$ está bien definida debemos comprobar que fijado $k \in K$, saca escalares, pues ya sabemos que es homomorfismo de grupos.
        \begin{equation*}
            h(k)(jv) = kjv = jkv = jh(k)(v) \quad \forall j \in K
        \end{equation*}
    \item Si definimos $h : K \to End(V)$, esto nos da la estructura de $K$-módulo, por lo que es un homomorfismo de anillos, y esto no se pierde al restringir el codominio. 
    \item Veamos la última afirmación del enunciado:
        \begin{equation*}
            (h(k)\circ T)(v) = h(k)(T(v)) = kT(v) = T(kv) = T(h(k)(v)) = (T \circ h(k))(v)
        \end{equation*}
        Por tanto, $h(k)$ conmuta con todo endomorfismo $K$-lineal, por lo cual el resto de afirmaciones ya son claras, con la estructura de $K$-álgebra definida por el homomorfismo $h$. 
\end{itemize}
\end{ejercicio} 

%--------------------------------------------------------------
\begin{ejercicio}
    Supongamos que $A$ y $B$ son $K$-álgebras con morfismos de estructura $\rho_A:K \to A$ y $\rho_B:K \to B$. Sea $\phi:A \to B$ un homomorfismo de anillos. Demostrar que $\phi $ es homomorfismo de $K$-álgebras si, y sólo si, $\phi \circ \rho_A = \rho_B$. \\

\noindent
\textbf{Solución:}
\begin{description}
    \item [$\Longrightarrow )$] $\phi(\rho_A(k)) = \phi(k1_A) \AstIg k\phi(1_A) = k1_B = \rho_B(k)$ $\forall k \in K$. 
    Donde en $(\ast)$ hemos usado que $\phi$ es $K$-lineal.
    \item [$\Longleftarrow )$] Nos falta ver que $\phi$ es $K$-lineal, en particular, que saca escalares. Sea $k \in K$ y $a \in A$:
        \begin{equation*}
            \phi(ka) = \phi(k1_Aa) = \phi(k1_A)\phi(a) = \phi(\rho_A(k))\phi(a) \AstIg \rho_B(k)\phi(a) = k1_B\phi(a) = k\phi(a),
        \end{equation*}
        donde en $(\ast)$ hemos usado la hipótesis.
\end{description}
\end{ejercicio}

%--------------------------------------------------------------
\begin{ejercicio}
    Sea $A$ un espacio vectorial sobre un cuerpo $K$, y fijemos esta estructura. La acción de un escalar de $K$ sobre un vector de $A$ se denotará por yuxtaposición. Demostrar que dar una estructura de $K$-álgebra asociativa unital sobre $A$ es equivalente a dar una multiplicación asociativa $\ast:A\times A \to A$ y $K$-bilineal junto con una aplicación $K$-lineal $\eta:K \to A$ tal que $\eta(k)\ast a = a \ast \eta(k)$ para todo $k \in K, a \in A$. \\
    
    \noindent
    \textbf{Solución:}
    
\end{ejercicio}
%--------------------------------------------------------------
\begin{ejercicio}
    Comprobar todas las afirmaciones hechas en el Ejemplo 1.11.
\end{ejercicio}

%--------------------------------------------------------------
\begin{ejercicio}
    Sea $K$ un cuerpo. Dar la lista, salvo isomorfismmos, de todas las $K$-álgebras asociativas unitales de dimensión 2.
\end{ejercicio}

%--------------------------------------------------------------
\begin{ejercicio}
    Expresar el cuerpo $\mathbb{Q}(\sqrt{2})$ como una $\mathbb{Q}$-subálgebra de un álgebra de matrices sobre $\mathbb{Q}$. \\
    
    \noindent
    \textbf{Solución:}
    Por conocimientos de Álgebra III, sabemos que una $\mathbb{Q}$-base de $\mathbb{Q}(\sqrt{2})$ es $\{1, \sqrt{2}\}$. Sea pues el homomorfismo de estructura de $\mathbb{Q}$-módulo $\lambda: \mathbb{Q}(\sqrt{2}) \to End(\mathbb{Q})$. Sea $a = \alpha + \beta \sqrt{2}$, veamos cuales son las imágenes de los elementos de la base por $\lambda(a)$:
    \begin{gather*}
        \lambda(a)(1) = a = \alpha + \beta \sqrt{2} \\
        \lambda(a)(\sqrt{2}) = \alpha\sqrt{2} + 2\beta 
    \end{gather*}
    Si ahora esto lo escribimos en coordenadas de la base que tenemos, obtenemos lo siguiente:
    \begin{equation*}
        M_{\cc{B}} = \left(\begin{array}{cc}
            \alpha & 2\beta \\
            \beta & \alpha 
        \end{array}\right)
    \end{equation*}
    Por tanto, obtenemos que:
    \begin{equation*}
        \mathbb{Q}(\sqrt{2}) \cong \left\{\left(\begin{array}{cc}
            \alpha & 2\beta \\
            \beta & \alpha 
        \end{array}\right) : \alpha, \beta \in \mathbb{Q}\right\} \subset M_2(\mathbb{Q})
    \end{equation*}
\end{ejercicio}

%--------------------------------------------------------------
\begin{ejercicio}
    Sea 
    \begin{equation*}
        \mathbb{H} = \left\{\left(\begin{array}{cc}
            \alpha & -\overline{\beta} \\
            \beta & \overline{\alpha } 
        \end{array}\right) : \alpha, \beta \in \mathbb{C}\right\}
    \end{equation*}
    \begin{enumerate}
        \item Demostrar que $\mathbb{H}$ es un subálgebra real de $M_2(\mathbb{C})$ y que $Z(\mathbb{H}) = \mathbb{R}$.
        \item Demostrar que todo elemento no nulo de $\mathbb{H}$ es una unidad.
        \item Demostrar que las matrices 
            TODO: meter aqui las matrices
            forma una base de $\bb{H}$ como espacio vectorial real.
            \item Comprobar las identidades: $i^2 = j^2 = k^2 = -1, ij=k, jk=i, ki=j$.
    \end{enumerate}
\end{ejercicio}

%--------------------------------------------------------------
\begin{ejercicio}
    Dado un $A$-módulo $V$ no nulo, demostrar que
    \begin{equation*}
        Ann_A(V) = \left\{a \in A : av = 0 \ \forall v \in V\right\}
    \end{equation*}
        es un ideal de $A$. Dotar a $V$ de estructura de $A \slash Ann_A(V)$-módulo fiel. \\
        
    \noindent
    \textbf{Solución:}
    Vamos a probar que es un ideal de $A$:
    \begin{description}
        \item [Cerrado para producto por elementos de $A$:] Sea $a, c \in A, m \in M$:
            \begin{gather*}
                (ca)m = c(am) = c0 = 0 \\
                (ac)m = a(cm) = 0
            \end{gather*}
            En el segundo caso $cm \in M$, por lo que tenemos el resultado.
        \item [Subgrupo aditivo de $A$:] Sean $a, b \in Ann_A(V)$:
            \begin{equation*}
                (a-b)m = am + (-b)m = 0 + (-1b)m = 0
            \end{equation*}
            La última igualdad se debe al punto anterior y que $-1 \in A$.
    \end{description}

    Ahora debemos dotar a $V$ de estructura de $A/Ann_A(V)$-módulo. Sabemos que es $A$-módulo, luego existe un homomorfismo de anillos: $\phi:A \to End(V)$ definido por $\phi(a)(v) = a\cdot v$, siendo $\cdot$ el producto de escalares. Bien, pues veamos que $Ann_A(V)= ker(\phi)$. 
    \begin{equation*}
        a \in ker(\phi) \Longleftrightarrow 0 = \phi(a)(v) = av \quad \forall v \in V \Longleftrightarrow a \in Ann_A(V)
    \end{equation*}
    Por tanto, por el Primer Teorema de Isomorfía para anillos:
    \begin{equation*}
        \phi:A/Ann_A(V) \to End(V)
    \end{equation*}
    es un homomorfismo de anillos inyectivo y hemos terminado. 
        
\end{ejercicio}

%--------------------------------------------------------------
\begin{ejercicio}
    Dar una demostración del Lema 1.4.2.
\end{ejercicio}

%--------------------------------------------------------------
\begin{ejercicio}
    
\end{ejercicio}
